% 2026-11-09 站主問本週三篇對財政法、稅捐法的意義與美國法上的聯想；底稿同日記入 notes/docs/1151/seminar/問答.md。
% 美國判決照既有體例行內引，不入 refs.bib。案號、卷頁、日期 2026-11-09 經 agent 核對（LII、聯邦憲法法院、憲法委員會、憲法法庭網站）；
% Loper Bright、Mayo、NFIB 之卷頁與 Golsen 僅經 Wikipedia 核；IRS nonacquiescence 之效果未取得 IRS 官方說明；BVerfGE 93, 121 主文未取原文。
\chapter{與財政法、稅捐法的關係}
\label{ch:55_fiscal}

\begin{qaquestion}
Alexander\ix{persons}{alexander@Alexander, Larry} 與 Schauer\ix{persons}{schauer@Schauer, Frederick}、Kramer\ix{persons}{kramer@Kramer, Larry D.}、Ferejohn\ix{persons}{ferejohn@Ferejohn, John} 與 Pasquino\ix{persons}{pasquino@Pasquino, Pasquale} 三篇，對財政法與稅捐法有何意義？美國法上有哪些可以對照的案件與制度？
\end{qaquestion}

三篇皆未以租稅為例；以下所引租稅案件與制度，均為三篇以外之對照。

\section{一、Alexander 與 Schauer：解釋權威與定紛止爭}
\label{sec:55_fiscal-1}

\begin{enumerate}
\tightlist
\item
  兩人主張其他部門應服從最高法院對憲法的解釋，理由在於權威性解釋所提供的定紛止爭（settlement\ix{terms}{02@\zhlidxcat{定紛止爭論}!定紛止爭}）利益：穩定、可預測與協調（457 至 463 頁）。批評者無人主張下級法院無須遵循最高法院，也無人支持 Meese\ix{persons}{meese@Meese, Edwin, III} 主張的各州得將最高法院裁判視為只拘束該案當事人的立場（476 頁）。\footcite[457--463, 476]{alexander2000}
\item
  美國聯邦稅務局（IRS）對其敗訴的稅務法院（Tax Court）判決，於 Internal Revenue Bulletin 公告 acquiescence（接受）或 nonacquiescence（不接受）；公告不接受者，該局在其他納稅義務人的案件中仍維持原有見解。稅務法院自身則依 Golsen v. Commissioner\ix{cases}{02@美國!03@稅務法院!1970Golsen v. Commissioner（1970）@Golsen v. Commissioner（1970）}, 54 T.C. 742 (1970) 遵循該案上訴所屬巡迴上訴法院的先例。兩項制度都以下級法院的法律解釋為對象，與 Cooper 規則\ix{terms}{01@\zhlidxcat{司法審查權限之歸屬}!Cooper 規則}所涉最高法院的憲法解釋不在同一層次；兩人所稱的定紛止爭利益，在租稅行政上須與全國課稅之一致性相權衡，兩人文中未處理此項權衡。
\item
  Loper Bright Enterprises v. Raimondo\ix{cases}{02@美國!01@聯邦最高法院!2024Loper Bright Enterprises v. Raimondo（2024）@Loper Bright Enterprises v. Raimondo（2024）}, 603 U.S. 369 (2024) 廢棄 Chevron 原則，法律之解釋改由法院獨立判斷，不再尊重\ix{terms}{05@\zhlidxcat{解釋理論與自制}!尊重}行政機關之合理解釋。Mayo Foundation for Medical Education \& Research v. United States, 562 U.S. 44 (2011) 曾明示 Chevron 原則適用於財政部之稅法規則。Loper Bright 所涉為法律解釋，兩人所論為憲法解釋；兩人對解釋權威歸屬的論證能否延伸至法律解釋，是可以提出的問題。
\item
  中華民國：釋字第 217 號\ix{cases}{01@中華民國!01@司法院大法官!0217@釋字第 217 號}解釋以憲法第 19 條\ix{statutes}{01@中華民國!01@中華民國憲法!019@第 19 條}\footcite[第 19 條]{roc-constitution}為據，界定人民僅依法律所定之納稅主體、稅目、稅率、納稅方法及納稅期間等項負納稅義務。\footcite{jy217}大法官解釋與憲法法庭判決拘束各機關及人民（釋字第 185 號\ix{cases}{01@中華民國!01@司法院大法官!0185@釋字第 185 號}解釋；憲法訴訟法第 38 條\ix{statutes}{01@中華民國!02@憲法訴訟法!038@第 38 條}第 1 項）。依兩人 457 頁的循環論證之說，此一拘束力以法律明定之後，是否仍須另以定紛止爭利益說明其正當性，可列為課前問題。
\end{enumerate}

\section{二、Kramer：自制的兩種理由與司法至上的興起}
\label{sec:55_fiscal-2}

\begin{enumerate}
\tightlist
\item
  Kramer\ix{persons}{kramer@Kramer, Larry D.} 626 至 630 頁敘述重建之後司法至上\ix{terms}{01@\zhlidxcat{司法審查權限之歸屬}!司法至上}強勢重現，舉 Brewer 大法官 1893 年的演說，並稱進步派與法院的衝突延續至新政危機，法院在 1930 年代對新政立法「tried, too late, to draw a line」（630 頁）。\footcite[626--630]{kramer2012}
\item
  Pollock\ix{cases}{02@美國!01@聯邦最高法院!1895Pollock v. Farmers' Loan \& Trust Co.（1895）@Pollock v. Farmers' Loan \& Trust Co.（1895）} v. Farmers' Loan \& Trust Co., 157 U.S. 429, 158 U.S. 601 (1895) 宣告 1894 年所得稅法違憲，理由為所得稅屬直接稅而未依各州人口分配。國會與各州其後以第十六修正案（1913 年生效）准許聯邦課徵所得稅而不依人口分配。Kramer 未引 Pollock，該案作成於 Kramer 所述 1893 年以後法院與進步派衝突的期間。推翻 Pollock 的途徑是修憲\ix{terms}{03@\zhlidxcat{憲法之效力基礎}!修憲}，各部門並未拒絕執行該判決；Alexander\ix{persons}{alexander@Alexander, Larry} 與 Schauer\ix{persons}{schauer@Schauer, Frederick} 466 頁所提第 5 條修憲程序，即此一途徑。
\item
  新政時期，United States v. Butler\ix{cases}{02@美國!01@聯邦最高法院!1936United States v. Butler（1936）@United States v. Butler（1936）}, 297 U.S. 1 (1936) 宣告農業調整法以加工稅支應之生產調整計畫違憲；Steward Machine Co.~v. Davis, 301 U.S. 548 (1937) 維持社會安全法的失業保險稅。兩案相隔一年，正落在 Kramer 630 頁所述法院對新政劃線而後退讓的時期。
\item
  National Federation of Independent Business v. Sebelius\ix{cases}{02@美國!01@聯邦最高法院!2012NFIB v. Sebelius（2012）@NFIB v. Sebelius（2012）}, 567 U.S. 519 (2012) 以課稅權維持健保法個人強制納保之規定，與 Kramer 一文同年。Kramer 633 頁說，在司法至上之下，自制\ix{terms}{05@\zhlidxcat{解釋理論與自制}!司法自制}只剩不致引發政治反撲的審慎\ix{terms}{05@\zhlidxcat{解釋理論與自制}!審慎}理由。該案以課稅權為依據而維持法律，屬於 Kramer 所說共和派\ix{terms}{01@\zhlidxcat{司法審查權限之歸屬}!共和派}以其他機關另有解釋權威為前提的自制，或屬審慎的權宜，可依 Kramer 的分類提問。
\end{enumerate}

\section{三、Ferejohn 與 Pasquino：審議型態與立法者角色}
\label{sec:55_fiscal-3}

\begin{enumerate}
\tightlist
\item
  兩人說明法國憲法委員會於法律公布前審查，延遲通常為一個月，政府宣告緊急時縮短為八日，審查程序承受時間壓力（33 頁）。\footcite[33]{ferejohn2002}財政法案另受第五共和憲法第 47 條所定議決期限之限制。
\item
  兩人歸納 Kelsen\ix{persons}{kelsen@Kelsen, Hans（凱爾森）} 的觀念：憲法法院宣告法律違憲，須重建該法律之前的法律狀態，因而兼有積極立法的作用（30 頁，註 13）。德國聯邦憲法法院於租稅案件有宣告法律與基本法不符而不宣告無效、並定期限命立法者修正之例，如財產稅裁判（BVerfGE 93, 121，1995 年）與不動產稅（Grundsteuer）裁判（BVerfGE 148, 147，2018 年）。租稅法律一經宣告無效即生稅收短缺，立法者須另定替代規定，法院之裁判因此兼及立法上的安排。
\end{enumerate}
